\section{Keeping Advent}

\subsection{Watch}\label{sec:watch}

For the time before his coming the Lord left one word, and he spoke it to
everyone:

\begin{quote}
Take ye heed, watch and pray. For ye know not when the time is. \dots\ Watch ye
therefore (for you know not when the lord of the house cometh, at even, or at
midnight, or at the cock crowing, or in the morning): Lest coming on a sudden,
he find you sleeping. And what I say to you, I say to all: Watch.\par
{\smallskip\noindent\small ---Mark 13:33--37\par}
\end{quote}

To watch is to stay awake for someone who is coming. It is not the same as
knowing when he will come. St Cyril of Jerusalem warned his hearers against
that curiosity: ``The things then which are seen
shall pass away, and there shall come the things which are looked for, things
fairer than the present; but as to the time let no one be curious.'' Watching
does not mean reckoning dates. It means living now as one who expects to be
visited.

\subsubsection*{Watching that does something}

St Cyril imagines someone asking how a man is to escape the fire and enter the
kingdom, and he does not reach for an allegory:

\begin{quote}
I was an hungred, He says, and ye gave Me meat. Learn hence the way; there is
here no need of allegory, but to fulfil what is said. I was an hungred, and ye
gave Me meat; I was thirsty, and ye gave Me drink; I was a stranger, and ye took
Me in; naked, and ye clothed Me; I was sick, and ye visited Me; I was in prison,
and ye came unto Me. These things if thou do, thou shalt reign together with
Him; but if thou do them not, thou shalt be condemned. At once then begin to do
these works, and abide in the faith; lest, like the foolish virgins, tarrying
to buy oil, thou be shut out. Be not confident because thou merely possessest
the lamp, but constantly keep it burning.\par
{\smallskip\noindent\small ---St Cyril of Jerusalem, Catechetical Lecture 15\par}
\end{quote}

The ten virgins in the parable all had lamps, and all of them slept while the
bridegroom tarried; the difference was the oil.

There is a promise hidden in another of the Lord's sayings about watching.
``Blessed are those servants whom the Lord, when he cometh, shall find
watching. Amen I say to you that he will gird himself and make them sit down to
meat and passing will minister unto them.'' The master comes home, finds his
servants awake, and serves them at table.

\subsubsection*{Watching that prays}

St Augustine, in a letter on prayer, says that desire is itself a kind of
prayer:

\begin{quote}
When we cherish uninterrupted desire along with the exercise of faith and hope
and charity, we ``pray always.'' But at certain stated hours and seasons we
also use words in prayer to God, that by these signs of things we may admonish
ourselves, and may acquaint ourselves with the measure of progress which we
have made in this desire, and may more warmly excite ourselves to obtain an
increase of its strength.\par
{\smallskip\noindent\small ---St Augustine, Letter 130\par}
\end{quote}

The words said at morning, noon and evening are not the whole of prayer; they
keep the desire awake between them. Augustine says in another place that if
your desire continues without ceasing, your prayer continues too.

The shepherds of Bethlehem were awake. They were ``watching and keeping the
night watches over their flock'', and so they were the first to hear. St
Gregory the Great saw a reason in it:
the angel appeared to them ``while they were watching'', because ``they are
thought worthy above the rest to see sublime things who take a watchful care
of their faithful flocks''. His words reach anyone who keeps watch over what is
given into his care: a family, a parish, a class of children, a sick friend.

Watching also receives the hidden coming of which St Bernard spoke, the Lord's
coming into men between the first coming and the last. The Lord promised it
Himself: ``If any one love me, he will keep my word. And my Father will love
him and we will come to him and will make our abode with him.''

Choose one of the works in Cyril's list and do it before Christmas: food for
someone hungry, a visit to someone sick or in prison, a welcome for a stranger.
Then choose the hours of the day at which you will use words in prayer, as
Augustine says, and keep them.
